\chapter{Deployment}

\begin{notebox}[NOTE]
Concepts and config values a production MERN deployment needs, provider-agnostic. Git workflow and provider CLI/sign-up steps are out of scope.
\end{notebox}

\section{Production Architecture}

\begin{diagrambox}[Production Request Path]
\begin{verbatim}
Internet
   |  HTTPS
   v
React Frontend (static hosting, e.g. CDN)
   |  HTTPS API calls (VITE_API_URL)
   v
Express Backend (Node hosting)
   |  MongoDB driver (TLS connection string)
   v
MongoDB Atlas (managed cluster)
\end{verbatim}
\end{diagrambox}

Two independent deployments connected only by URLs and env vars: frontend is a static build artifact, backend is a long-running Node process, database is a separate managed service.

\section{Frontend: Build for Production}

\begin{lstlisting}
npm run build
# produces a dist/ folder: static HTML/CSS/JS, ready for any static host
\end{lstlisting}

Static host serves \texttt{dist/} directly --- no Node process for the frontend. Set \texttt{VITE\_API\_URL} to the backend's URL \emph{before} running \texttt{build}; Vite inlines env vars at build time, not runtime.

\begin{warnbox}[WARNING]
If \texttt{VITE\_API\_URL} still points at \texttt{localhost} at build time, every production API call targets the visitor's own machine. Rebuild after changing the value --- editing \texttt{.env} alone won't update an existing \texttt{dist/}.
\end{warnbox}

\section{Backend: Start Command}

\begin{lstlisting}
npm start
# package.json: "start": "node src/server.js"
\end{lstlisting}

Production runs \texttt{node}, never \texttt{nodemon} --- nodemon is for local file-watching only. Restarting on crash is the hosting platform's job.

\section{Environment Variables}

Configure these on the hosting platform's dashboard --- never commit \texttt{.env}.

\begin{tabularx}{\textwidth}{@{}>{\bfseries}p{4.2cm}X@{}}
\toprule
Variable & Production value \\
\midrule
\texttt{MONGO\_URI} & Atlas connection string for the production cluster \\
\texttt{JWT\_SECRET} & Long random string, distinct from dev secret \\
\texttt{NODE\_ENV} & \texttt{production} \\
\texttt{FRONTEND\_URL} & Deployed frontend's exact domain (for CORS) \\
\texttt{GOOGLE\_CLIENT\_ID/SECRET} & Production OAuth creds with redirect URI registered \\
\texttt{CLOUDINARY\_*} & Cloudinary account keys \\
\bottomrule
\end{tabularx}

\section{Updating CORS for Production}

\begin{lstlisting}[language=JavaScript]
app.use(cors({
  origin: process.env.FRONTEND_URL, // the deployed frontend's real domain
  credentials: true,
}));
\end{lstlisting}

A \texttt{localhost:5173} origin works locally but must be replaced with the deployed frontend's domain, or cookie-based requests are blocked in production (Chapter 26).

\section{MongoDB Atlas Network Access}

Atlas rejects IPs not on its allowlist. For dynamic outbound IPs, allow \texttt{0.0.0.0/0} (still gated by DB credentials). If the platform gives static IPs, allowlist those instead.

\section{HTTPS and Custom Domains}

Most platforms provision HTTPS automatically once a domain is attached. Once active, set cookies to \texttt{secure: true} and \texttt{sameSite: 'none'} for cross-domain delivery (Chapter 19). A custom domain is just a DNS record pointed at the platform.

\section{Debugging in Production}

Every platform exposes a logs view --- the production equivalent of local terminal output.

\begin{itemize}
\item Reproduce the request, check the platform's log viewer for the stack trace.
\item Rule out a missing/incorrect env variable (the most common production-only bug).
\item Check the deployed frontend's console and Network tab as you would locally.
\end{itemize}

\whatwhyhow
{Two independently deployed pieces --- static frontend build and long-running Node backend --- connected via env vars (\texttt{VITE\_API\_URL}, \texttt{FRONTEND\_URL}, \texttt{MONGO\_URI}) and a managed Atlas cluster.}
{Separating concerns lets the frontend scale as CDN static assets while the backend scales as a process, and keeps secrets out of the repo.}
{Build the frontend with the production API URL baked in, start the backend with \texttt{npm start}, configure CORS/cookies for real domains, set secrets via the platform's env system.}
